\documentclass{beamer}

\usetheme{Warsaw}
\usepackage{graphicx}
\usepackage{ulem}
\usepackage{tikz}
% include.tex

% \newcommand{\expit}[1]{\text{expit} #1}
% \newcommand{\logit}[1]{\text{logit} #1}

\def \R {{\mathbb{R}}}
\def \vbeta {{\boldsymbol \beta}}
\def \vnu {{\bf \nu}}
\def \vy {{\bf y}}
\def \vx {{\bf x}}
\def \vu {{\bf u}}
\def \vr {{\bf r}}
\def \vp {{\bf p}}
\def\vectorfontone{\bf}
\def\vone{{\bf 1}}
\def\vzero{{\bf 0}}
\def \vmu {{\boldsymbol \mu}}
\def \vnu {{\bf \nu}}
\def \vmuqbeta {{\vmu_{q(\vbeta)}}}
\def \vmubeta {{\vmu_{\vbeta}}}
\def \Sigmaqbeta {{\Sigma_{q(\vbeta)}}}
\def \Sigmabeta {{\Sigma_{\vbeta}}}
\def \va {{\bf a}}
\def \vtheta {{\bf \theta}}
\def \mX {{\bf X}}
\def \mZ {{\bf Z}}
\def \mR {{\bf R}}
\def \mC {{\bf C}}
\def \mI {{\bf I}}
\def \mLambda {{\boldsymbol \Lambda}}
\def \mSigma {{\boldsymbol \Sigma}}
\def \B {{\text{B}}}

\def\ds{{\displaystyle}}

\def\diag{{\mbox{diag}}}
\def\bbE{\mathbb{E}}


\usepackage{latexsym,amssymb,amsmath,amsfonts}
%\usepackage{tabularx}
\usepackage{theorem}
\usepackage{verbatim,array,multicol,palatino}
\usepackage{graphicx}
\usepackage{graphics}
\usepackage{fancyhdr}
\usepackage{algorithm,algorithmic}
\usepackage{url}
%\usepackage[all]{xy}



\def\approxdist{\stackrel{{\tiny \mbox{approx.}}}{\sim}}
\def\smhalf{\textstyle{\frac{1}{2}}}
\def\vxnew{\vx_{\mbox{{\tiny new}}}}
\def\bib{\vskip12pt\par\noindent\hangindent=1 true cm\hangafter=1}
\def\jump{\vskip3mm\noindent}
\def\etal{{\em et al.}}
\def\etahat{{\widehat\eta}}
\def\thick#1{\hbox{\rlap{$#1$}\kern0.25pt\rlap{$#1$}\kern0.25pt$#1$}}
\def\smbbeta{{\thick{\scriptstyle{\beta}}}}
\def\smbtheta{{\thick{\scriptstyle{\theta}}}}
\def\smbu{{\thick{\scriptstyle{\rm u}}}}
\def\smbzero{{\thick{\scriptstyle{0}}}}
\def\boxit#1{\begin{center}\fbox{#1}\end{center}}
\def\lboxit#1{\vbox{\hrule\hbox{\vrule\kern6pt
      \vbox{\kern6pt#1\kern6pt}\kern6pt\vrule}\hrule}}
\def\thickboxit#1{\vbox{{\hrule height 1mm}\hbox{{\vrule width 1mm}\kern6pt
          \vbox{\kern6pt#1\kern6pt}\kern6pt{\vrule width 1mm}}
               {\hrule height 1mm}}}


%\sloppy
%\usepackage{geometry}
%\geometry{verbose,a4paper,tmargin=20mm,bmargin=20mm,lmargin=40mm,rmargin=20mm}


%%%%%%%%%%%%%%%%%%%%%%%%%%%%%%%%%%%%%%%%%%%%%%%%%%%%%%%%%%%%%%%%%%%%%%%%%%%%%%%%
%
% Some convenience definitions
%
% \bf      -> vector
% \sf      -> matrix
% \mathcal -> sets or statistical
% \mathbb  -> fields or statistical
%
%%%%%%%%%%%%%%%%%%%%%%%%%%%%%%%%%%%%%%%%%%%%%%%%%%%%%%%%%%%%%%%%%%%%%%%%%%%%%%%%

% Sets or statistical values
\def\sI{{\mathcal I}}                            % Current Index set
\def\sJ{{\mathcal J}}                            % Select Index set
\def\sL{{\mathcal L}}                            % Likelihood
\def\sl{{\ell}}                                  % Log-likelihood
\def\sN{{\mathcal N}}                            
\def\sS{{\mathcal S}}                            
\def\sP{{\mathcal P}}                            
\def\sQ{{\mathcal Q}}                            
\def\sB{{\mathcal B}}                            
\def\sD{{\mathcal D}}                            
\def\sT{{\mathcal T}}
\def\sE{{\mathcal E}}                            
\def\sF{{\mathcal F}}                            
\def\sC{{\mathcal C}}                            
\def\sO{{\mathcal O}}                            
\def\sH{{\mathcal H}} 
\def\sR{{\mathcal R}}                            
\def\sJ{{\mathcal J}}                            
\def\sCP{{\mathcal CP}}                            
\def\sX{{\mathcal X}}                            
\def\sA{{\mathcal A}} 
\def\sZ{{\mathcal Z}}                            
\def\sM{{\mathcal M}}                            
\def\sK{{\mathcal K}}     
\def\sG{{\mathcal G}}                         
\def\sY{{\mathcal Y}}                         
\def\sU{{\mathcal U}}  


\def\sIG{{\mathcal IG}}                            


\def\cD{{\sf D}}
\def\cH{{\sf H}}
\def\cI{{\sf I}}

% Vectors
\def\vectorfontone{\bf}
\def\vectorfonttwo{\boldsymbol}
\def\va{{\vectorfontone a}}                      %
\def\vb{{\vectorfontone b}}                      %
\def\vc{{\vectorfontone c}}                      %
\def\vd{{\vectorfontone d}}                      %
\def\ve{{\vectorfontone e}}                      %
\def\vf{{\vectorfontone f}}                      %
\def\vg{{\vectorfontone g}}                      %
\def\vh{{\vectorfontone h}}                      %
\def\vi{{\vectorfontone i}}                      %
\def\vj{{\vectorfontone j}}                      %
\def\vk{{\vectorfontone k}}                      %
\def\vl{{\vectorfontone l}}                      %
\def\vm{{\vectorfontone m}}                      % number of basis functions
\def\vn{{\vectorfontone n}}                      % number of training samples
\def\vo{{\vectorfontone o}}                      %
\def\vp{{\vectorfontone p}}                      % number of unpenalized coefficients
\def\vq{{\vectorfontone q}}                      % number of penalized coefficients
\def\vr{{\vectorfontone r}}                      %
\def\vs{{\vectorfontone s}}                      %
\def\vt{{\vectorfontone t}}                      %
\def\vu{{\vectorfontone u}}                      % Penalized coefficients
\def\vv{{\vectorfontone v}}                      %
\def\vw{{\vectorfontone w}}                      %
\def\vx{{\vectorfontone x}}                      % Covariates/Predictors
\def\vy{{\vectorfontone y}}                      % Targets/Labels
\def\vz{{\vectorfontone z}}                      %

\def\vone{{\vectorfontone 1}}
\def\vzero{{\vectorfontone 0}}

\def\valpha{{\vectorfonttwo \alpha}}             %
\def\vbeta{{\vectorfonttwo \beta}}               % Unpenalized coefficients
\def\vgamma{{\vectorfonttwo \gamma}}             %
\def\vdelta{{\vectorfonttwo \delta}}             %
\def\vepsilon{{\vectorfonttwo \epsilon}}         %
\def\vvarepsilon{{\vectorfonttwo \varepsilon}}   % Vector of errors
\def\vzeta{{\vectorfonttwo \zeta}}               %
\def\veta{{\vectorfonttwo \eta}}                 % Vector of natural parameters
\def\vtheta{{\vectorfonttwo \theta}}             % Vector of combined coefficients
\def\vvartheta{{\vectorfonttwo \vartheta}}       %
\def\viota{{\vectorfonttwo \iota}}               %
\def\vkappa{{\vectorfonttwo \kappa}}             %
\def\vlambda{{\vectorfonttwo \lambda}}           % Vector of smoothing parameters
\def\vmu{{\vectorfonttwo \mu}}                   % Vector of means
\def\vnu{{\vectorfonttwo \nu}}                   %
\def\vxi{{\vectorfonttwo \xi}}                   %
\def\vpi{{\vectorfonttwo \pi}}                   %
\def\vvarpi{{\vectorfonttwo \varpi}}             %
\def\vrho{{\vectorfonttwo \rho}}                 %
\def\vvarrho{{\vectorfonttwo \varrho}}           %
\def\vsigma{{\vectorfonttwo \sigma}}             %
\def\vvarsigma{{\vectorfonttwo \varsigma}}       %
\def\vtau{{\vectorfonttwo \tau}}                 %
\def\vupsilon{{\vectorfonttwo \upsilon}}         %
\def\vphi{{\vectorfonttwo \phi}}                 %
\def\vvarphi{{\vectorfonttwo \varphi}}           %
\def\vchi{{\vectorfonttwo \chi}}                 %
\def\vpsi{{\vectorfonttwo \psi}}                 %
\def\vomega{{\vectorfonttwo \omega}}             %


% Matrices
%\def\matrixfontone{\sf}
%\def\matrixfonttwo{\sf}
\def\matrixfontone{\bf}
\def\matrixfonttwo{\boldsymbol}
\def\mA{{\matrixfontone A}}                      %
\def\mB{{\matrixfontone B}}                      %
\def\mC{{\matrixfontone C}}                      % Combined Design Matrix
\def\mD{{\matrixfontone D}}                      % Penalty Matrix for \vu_J
\def\mE{{\matrixfontone E}}                      %
\def\mF{{\matrixfontone F}}                      %
\def\mG{{\matrixfontone G}}                      % Penalty Matrix for \vu
\def\mH{{\matrixfontone H}}                      %
\def\mI{{\matrixfontone I}}                      % Identity Matrix
\def\mJ{{\matrixfontone J}}                      %
\def\mK{{\matrixfontone K}}                      %
\def\mL{{\matrixfontone L}}                      % Lower bound
\def\mM{{\matrixfontone M}}                      %
\def\mN{{\matrixfontone N}}                      %
\def\mO{{\matrixfontone O}}                      %
\def\mP{{\matrixfontone P}}                      %
\def\mQ{{\matrixfontone Q}}                      %
\def\mR{{\matrixfontone R}}                      %
\def\mS{{\matrixfontone S}}                      %
\def\mT{{\matrixfontone T}}                      %
\def\mU{{\matrixfontone U}}                      % Upper bound
\def\mV{{\matrixfontone V}}                      %
\def\mW{{\matrixfontone W}}                      % Variance Matrix i.e. diag(b'')
\def\mX{{\matrixfontone X}}                      % Unpenalized Design Matrix/Nullspace Matrix
\def\mY{{\matrixfontone Y}}                      %
\def\mZ{{\matrixfontone Z}}                      % Penalized Design Matrix/Kernel Space Matrix

\def\mGamma{{\matrixfonttwo \Gamma}}             %
\def\mDelta{{\matrixfonttwo \Delta}}             %
\def\mTheta{{\matrixfonttwo \Theta}}             %
\def\mLambda{{\matrixfonttwo \Lambda}}           % Penalty Matrix for \vnu
\def\mXi{{\matrixfonttwo \Xi}}                   %
\def\mPi{{\matrixfonttwo \Pi}}                   %
\def\mSigma{{\matrixfonttwo \Sigma}}             %
\def\mUpsilon{{\matrixfonttwo \Upsilon}}         %
\def\mPhi{{\matrixfonttwo \Phi}}                 %
\def\mOmega{{\matrixfonttwo \Omega}}             %
\def\mPsi{{\matrixfonttwo \Psi}}                 %

\def\mone{{\matrixfontone 1}}
\def\mzero{{\matrixfontone 0}}

% Fields or Statistical
\def\bE{{\mathbb E}}                             % Expectation
\def\bP{{\mathbb P}}                             % Probability
\def\bR{{\mathbb R}}                             % Reals
\def\bI{{\mathbb I}}                             % Reals
\def\bV{{\mathbb V}}                             % Reals

\def\vX{{\vectorfontone X}}                      % Targets/Labels
\def\vY{{\vectorfontone Y}}                      % Targets/Labels
\def\vZ{{\vectorfontone Z}}                      %

% Other
\def\etal{{\em et al.}}
\def\ds{\displaystyle}
\def\d{\partial}
\def\diag{\text{diag}}
%\def\span{\text{span}}
\def\blockdiag{\text{blockdiag}}
\def\tr{\text{tr}}
\def\RSS{\text{RSS}}
\def\df{\text{df}}
\def\GCV{\text{GCV}}
\def\AIC{\text{AIC}}
\def\MLC{\text{MLC}}
\def\mAIC{\text{mAIC}}
\def\cAIC{\text{cAIC}}
\def\rank{\text{rank}}
\def\MASE{\text{MASE}}
\def\SMSE{\text{SASE}}
\def\sign{\text{sign}}
\def\card{\text{card}}
\def\notexp{\text{notexp}}
\def\ASE{\text{ASE}}
\def\ML{\text{ML}}
\def\nullity{\text{nullity}}

\def\logexpit{\text{logexpit}}
\def\logit{\mbox{logit}}
\def\dg{\mbox{dg}}

\def\Bern{\mbox{Bernoulli}}
\def\sBernoulli{\mbox{Bernoulli}}
\def\sGamma{\mbox{Gamma}}
\def\sInvN{\mbox{Inv}\sN}
\def\sNegBin{\sN\sB}

\def\dGamma{\mbox{Gamma}}
\def\dInvGam{\mbox{Inv}\Gamma}

\def\Cov{\mbox{Cov}}
\def\Mgf{\mbox{Mgf}}

\def\mis{{mis}} 
\def\obs{{obs}}

\def\argmax{\operatornamewithlimits{\text{argmax}}}
\def\argmin{\operatornamewithlimits{\text{argmin}}}
\def\argsup{\operatornamewithlimits{\text{argsup}}}
\def\arginf{\operatornamewithlimits{\text{arginf}}}


\def\minimize{\operatornamewithlimits{\text{minimize}}}
\def\maximize{\operatornamewithlimits{\text{maximize}}}
\def\suchthat{\text{such that}}


\def\relstack#1#2{\mathop{#1}\limits_{#2}}
\def\sfrac#1#2{{\textstyle{\frac{#1}{#2}}}}


\def\comment#1{
\vspace{0.5cm}
\noindent \begin{tabular}{|p{14cm}|}  
\hline #1 \\ 
\hline 
\end{tabular}
\vspace{0.5cm}
}


\def\mytext#1{\begin{tabular}{p{13cm}}#1\end{tabular}}
\def\mytextB#1{\begin{tabular}{p{7.5cm}}#1\end{tabular}}
\def\mytextC#1{\begin{tabular}{p{12cm}}#1\end{tabular}}

\def\jump{\vskip3mm\noindent}

\def\KL{\text{KL}}
\def\N{\text{N}}
\def\Var{\text{Var}}

\def \E {\mathbb{E}}
\def \BigO {\text{O}}
\def \IG {\text{IG}}
\def \Beta {\text{Beta}}



\usefonttheme{serif}

\title{Progress update}
\author{Mark Greenaway\\PhD candidate\\markg@maths.usyd.edu.au}

\mode<presentation>
{ \usetheme{boxes} }

\begin{document}
\begin{frame}
\titlepage
\end{frame}

\begin{frame}{Melbourne trip}
\begin{itemize}
\item Met Hadley Wickham. The future of R is in good hands. \\
	When I told him what I was doing, he responded ``I don't care about matrices.''. He only really cares 
	about data frames. \\
	A talk similiar to the talk he gave is available here \url{https://www.youtube.com/watch?v=hRNUgwAFZtQ}.
\item Hung out with Charles Gray, PhD candidate at Latrobe working with Agus. I taught her some UNIX and Git.
	She taught me some TikZ and LaTeX.\\
	\begin{tikzpicture}[scale=0.5]
	\draw (-0.75, 0.5) circle(10pt);
	\draw (0.75, 0.5) circle(10pt);
	\draw (-1, -0.75) to[out=-90, in=-90] (1, -0.75);
	\end{tikzpicture}
\item Met with Lev Lafayette. Sys admin. Set up USyd HPC (1,500 cores), now working with NeCTAR (4,000 cores). 
	He's a nice guy, and we had a productive chat.
\end{itemize}
\end{frame}

\def\vGamma{{\boldsymbol{\Gamma}}}
\begin{frame}{Focus this week: VB approximation for linear model selection}

Consider the model $\vy | \vbeta, \sigma^2 \sim \text{N}_n(\mX \vGamma \vbeta, \sigma^2 \mI)$ 
with priors
\begin{equation*}
\begin{array}{ll}
\vbeta | \sigma^2, g &\sim \text{N}_p(\vzero, g \sigma^2(\mX^\top \vGamma \mX)^{-1}), \\
\vgamma_i &\sim \text{Bernoulli}(\rho), 1 \leq i \leq n, 0 \leq \rho \leq 1, \\
p(\sigma^2) &= (\sigma^2)^{-1} I(\sigma^2 > 0) \text{ and }\\
p(g) &= \frac{g^b (1 + g)^{-a - b - 2}}{\text{Beta}(a + 1, b + 1)} I(g > 0)
\end{array}
\end{equation*}

where $\vGamma = \diag (\vgamma)$. Assume an approximation of the form

\begin{equation*}
q(\vtheta) = q(\vbeta) q(\vgamma) q(\sigma^2) q(g)
\end{equation*}
where
\begin{equation*}
\begin{array}{ll}
q^*(\vbeta) &= \text{N}(\vmu_{q(\vbeta), \mSigma_{q(\vbeta)}}), \\
q^*(\vgamma) &= \text{Bernoulli}(p_{q(\vgamma)}), \\
q^*(\sigma^2) &= \text{IG}(\alpha_{q(\sigma^2)}, \beta_{q(\sigma^2)}) \text{ and } \\
q^*(g) &= \text{Beta Prime}(\alpha_{q(g)}, \beta_{q(g)})). \\
\end{array}
\end{equation*}
\end{frame}

\begin{frame}{Mean field updates are coming}
\begin{equation*}
\begin{array}{rl}
&\quad \log{p(\vy, \vtheta)} \\
=& \quad \log{p(\vy | \vbeta, \vgamma, \sigma^2)} + \log{p(\vbeta | \vgamma, \sigma^2, g)} + \log{p(\sigma^2)} + \log{p(g)} \\
=& -\frac{n}{2} \log{2 \pi} -\half \log |\sigma^2 \mI_n| - \half (\vy - \mX \vGamma \vbeta)^\top \sigma^{-2} \mI_n (\vy - \mX \vGamma \vbeta) \\
& -\frac{p}{2} \log{2 \pi} - \half \log |g \sigma^2 (\mX^\top \vGamma \mX)^{-1}| -
	\half \vbeta^\top (g \sigma^2)^{-1} (\mX^\top \mX) \vbeta \\
& + \vone^\top \vgamma \log(\rho) + \vone^\top (\vone - \vgamma) \log(1 - \rho)\\
& - \log{(\sigma^2)} + \log \bI (\sigma^2 > 0) \\

\end{array}
\end{equation*}

Find the optimal mean field updates with
\[
q_{\vtheta_i} \propto \bE_{-\theta_i} [\log p(\vy, \vtheta)].
\]

\end{frame}

\begin{frame}{Algebra montage}

\begin{itemize}
\item I'm rusty on some things, and need to get my skills back up to scratch, particularly some linear 
algebra and multivariate distribution theory, as used in a Bayesian/VB context.
\item The only way to get better is to practice!
\begin{quote}
\tiny
The day's approaching to give it your best, \\
and you've got to reach your prime \\
That's when you need to put yourself to the test, \\
and show us a passage of time \\
We're gonna need a montage! (MONTAGE !) \\
A sports training montage! (MONTAGE !) \\
Show a lot of things happening at once, \\
remind everyone of what's going on! (WHAT'S GOING ON!) \\
And with every shot show a little improvement, \\
to show it all would take too long! \\
That's called a montage! (MONTAGE !) \\
Even Rocky had a montage! (MONTAGE !) \\
In any sport if you want to go, \\
from just a beginner to a pro \\
You need a montage! (MONTAGE !) \\
A simple little montage! (MONTAGE !) \\
Always fade out in a montage \\
If you fade out it seems like more time has a passed in a montage
\end{quote}

\end{itemize}
\end{frame}

\begin{frame}{Optimisation: Making numerical software run faster}
\begin{itemize}
\item The game is: 
\begin{enumerate}
\item Make sure your results are correct first
\item Measure/profile (don't guess!) to find out where time is being spent in your program
\item Make a change
\item Check your software still produces correct results and measure performance again. If it's slower, or the result is wrong, revert the change
\item Stop when your software is fast enough
\end{enumerate}
\item Don't skip steps!
\item I used \texttt{valgrind}.
\end{itemize}
\end{frame}

\begin{frame}{So let's play \ldots}
\begin{tabular}{|p{9cm}l|}
\hline
Change & Time (s) \\
\hline
Unoptimised code which produces correct results & 30 \\
Save common sub-expressions in rank-1 update & 18 \\
Use references to avoid copying & 16 \\
Pre-allocate memory for $\mX_\vgamma$ and $(\mX_\vgamma^\top \mX_\vgamma)^{-1}$ & 9 \\
\hline
\end{tabular}
\begin{itemize}
\item 0.472 seconds on verona using all cores. Over a million models a second
\item Other tricks I could try, but for now, that's fast enough!
\item I parallelised the code using OpenMP. Five lines.
\item I can recommend ``Using OpenMP'' by Barbara Chapman \note{MIT}, Gabrielle Jost, and Ruud Van der Pas.
\begin{figure}
\includegraphics[scale=0.3]{Using_OpenMP.jpeg}
\end{figure}
\note{The talks on Youtube by Tim Mattson (Intel) are also excellent.}
\end{itemize}
\end{frame}

\end{document}